% s-059-p059-m002; original PDF 59
\Aphorism{s-059-p059-m002}{56}
Ich kenne Menschen, die, wenn sie einen Andern stolpern oder springen sehen, am eigenen Leibe den Ruck empfinden, Menschen, denen es in Därmen und Nieren zuckt und zwackt, wenn sie von einer chirurgischen Operation lesen. Das „Mitleid“ ist hier zur physiologischen Wirklichkeit geworden, als eine Art gegenseitiger Induction der verschiedenen Nervensysteme unter sinnlicher oder gehirnlicher Vermittlung; und insofern jede Mitleidsmoral diese Inducibilität zu schärfen und zu erweitern trachtet, strebt sie einem unerträglichen Communismus des Leidens zu, wo die Individuation keine Schranke und die Haut kein Schutzfell mehr ist, wo Jedem seines „Bruders“ Auge und Zahn mit wehthut und jede Uebelkeit collectiv empfunden wird. Hofft man vielleicht, wie bei einem elektrischen Conductor, durch Vergrösserung der Oberfläche die „Spannung“ des Schmerzes zu mindern? Jedenfalls wird aus diesem künftigen Zustand heraus der Egoismus einst als Tugend und fast unerreichbares Ideal geflaggt werden, und jede Art Isolator wird hoch im Preise stehen — heute hat man noch seine unbedachte Freude, wenn von Socialorganismen, Volksseele oder Willensübertragung die Rede ist und dem Individuum die verfängliche Gelegenheit geboten wird, aus der Haut zu fahren.
\par

